\section{Introduction}
\label{sec:intro}

\subsection{The problem}

A workspace is a repository holding many independently versioned packages that
depend on one another. Releasing it means choosing a new version for each
package that needs one, writing those versions into manifests, uploading
artefacts to registries, and recording what happened.

The recording is the hard part. Registries have no transactions. A publish is an
irreversible write to a service the release tool does not control, several
registries may be involved (npm, a container registry, a language-specific
index), and there is no coordinator that can commit or abort them as a unit. A
run that publishes seven packages is seven independent, individually fallible
writes; it may finish having done four of them.

CCME\footnote{\emph{Conventional Commits: Monorepo Extension}, version
3.1.0-rc.2 at commit \texttt{252216ba}, \code{specs/ccme-spec/SPEC.md} in the
\code{dispat} repository. Referred to throughout as \SPEC{}, cited by section
number.} answers this the way a saga does. There is no external state file, no
lock service beyond the release lock itself, and no registry polling. The
release plan is a pure function of the repository at a fixed commit; each
completed leg of the transaction is recorded by an annotated git tag written
immediately after that leg succeeds; and recovery is re-execution: run the
engine again and it recomputes the same transaction, skipping the legs already
recorded.

That design only works if re-execution is well behaved, and \spec{13.7c} states
the eight properties that say it is:

\begin{description}[leftmargin=3.2em,style=nextline,font=\normalfont\bfseries]
  \item[\G{1} termination] propagation halts, whatever the declared depth.
  \item[\G{2} completeness] a contribution admitted to a dependent survives in
    every later run until that dependent releases at or after a release of
    every owed source carrying it; there are no orphans, whether the provider
    failed or the consumer got ahead of it.
  \item[\G{3} version stability] a package released late is released at the
    version it was originally planned at.
  \item[\G{4} no double ledger delivery] a contribution does not survive its
    delivery in the tag ledger; a consumer that proceeded past a commit before
    its source released it is bumped once more when that release comes, and
    that is the delivery, not a second one. This says nothing unconditional
    about external effects.
  \item[\G{5} no blast-radius widening] under the retry invariant, later
    targets are a subset of the first run's.
  \item[\G{6} restricted convergence] under the retry invariant, repeated
    successful progress empties the initial plan within as many runs as it had
    publishable packages.
  \item[\G{7} channel discharge] a prerelease channel is proposed to a package
    only until it is released on it.
  \item[\G{8} one-way axis independence] suppressing a version bump does not
    suppress a channel change; resolved channels may suppress bumps.
\end{description}

\subsection{Why formalise them, and what changed}

The specification proves each property briefly, while the algorithm is spread
across prose and pseudocode. This report writes the algorithm down once,
formally (\cref{sec:model,sec:propagation,sec:protocol}), and proves the
guarantees inside that model (\cref{sec:guarantees}), along with the
determinism and idempotency obligations of \spec{17.2}
(\cref{sec:determinism}), the duality of \spec{13.7b}
(\cref{sec:duality}), and the cost bounds and optimisation-equivalence claims
of \specs{13.11}{20.7} (\cref{sec:complexity}).

Earlier editions of this report formalised CCME 2.0.0, whose contribution was
to state the cross-run retry invariant explicitly and to separate the train
window from the undelivered window. The current text keeps both and changes
the admission test of the bump axis: where 2.0.0 admitted a propagated
contribution while the unit's commit lay in the dependent's undelivered
window, the candidate admits it while some source of the unit has not
\emph{delivered} that commit to the dependent, that is, while the dependent has
not released at or after a release of that source carrying it (\spec{13.4a}).
The two tests agree for every dependent that has not released past the commit.
They differ for a dependent that released past it before its provider did,
which the candidate's own publish protocol now allows: a consumer with a cause
of its own proceeds past a failed provider (\spec{19.3}) instead of being
blocked. Under the old test such a consumer was never planned again once the
provider published, an orphan reached from the consumer's side. Delivery
closes it, and two further rules follow from it: a window per dependency edge
that keeps the debt visible once the provider has released (\spec{13.3}), and a
refusal, \diag{E201}, where a provider would release at the very commit a
consumer it still owes was released at, because two releases on one commit
cannot be ordered afterwards. \Cref{sec:model,sec:propagation,sec:guarantees}
transcribe and prove the candidate's form; \cref{ob:samecommit} states the one
place where the ledger reports a loss instead of preventing it.

Around the computation the candidate adds what an engine that runs builds and
publishes across machines needs, and \cref{sec:protocol,sec:execution}
transcribe it as definitions and propositions rather than as narrative:
release records compared with the authoritative store under the lock before
planning (\spec{13.2}); a create-only rule for release tags (\spec{19.1}); the
three-valued build readiness relation of an edge, \code{none}, \code{build} or
\code{publish}, which is execution policy and never changes a plan
(\spec{19.2a}), with the run durations each value leads to; the distributed
execution profile (\spec{28}), where the plan is fixed under the lock before
any task is placed, outputs travel as verified first-class inputs, and
publication and recording stay serialised per owner; and the repository-local
windows of the polyrepository profile (\spec{27}). The single-history theorems
are not extended to several histories here; the report says so where the
objects are defined.

The reference implementation, dispat, implements the delivery admission and
the reconciliation of a proceeding consumer, and does not yet implement the
owed windows and \diag{E201}. \Cref{sec:obligations} keeps two ledgers: the
current departures of the engine from the candidate, each with the scenario
that exposes it, and the six failure classes found in CCME 1.0.0 with their
minimal counterexamples, which 2.0.0 repaired and which remain regression
evidence. \Cref{sec:conformance} maps each result to the specification's
conformance vectors and to the integration test that stands behind it.

\Cref{sec:evaluation} is new to this edition. It reports every executed
campaign with its pinned binary and environment: the differential validations
of the planner, the recovery and mixed-registry campaigns, the controlled fault
matrix with its finite projection, the planning-cost cells, and the two
measurements of 22 September 2026 on which the execution profile rests, the
Linux kernel farm groups and the LLVM matrix, each once, with the kit that
reproduces them.

\subsection{Scope, and what this is not}

The subject is the release computation of \spec{13}, the publish protocol of
\spec{19} and, as transcribed obligations, the execution rules of
\specs{19.2a}{28}; the guarantees are properties of the computation and the
protocol together, and no separation of them is meaningful. Message syntax
(\spec{4}--\spec{8}) is taken as given: units arrive already parsed, with their
scope-sets resolved and their corrections applied, which is what
\specs{13.4}{13.4b} deliver to \spec{13.5} onwards. The one exception is
\cref{prop:parse}, which states the parser's linearity because \spec{18.3}
depends on it.

Five things a complete account would also cover are outside the theorems, and
no theorem below should be read as saying anything about them: manifest range
reconciliation (\spec{9.4}, \spec{19.5}), which happens at publish time and
feeds back into the plan only through the published versions the protocol
names; the diagnostic registry of \spec{16}, of which only the codes naming a
suppression or a refusal are mentioned, as labels; the scheduling of the
executor beyond the happens-before relations and the duration bounds
\cref{sec:execution} states; shared-version groups (\spec{13.9a}), whose
member target floor and independent counter and channel axes are a
specification contract this report does not prove; and the adapter and
rollback protocols of \specs{25}{26}, which the current text marks as pending
implementation.

The report makes no novelty claim for the underlying patterns: sagas,
idempotent consumers, breadth-first search and list scheduling. Its
contribution is the formal model, corrected hypotheses, counterexamples,
conformance obligations, and the record of what was measured.

\subsection{Implementation examples outside the model}

Two small implementation examples show how the executor can be composed with
other controls; neither extends the theorem scope above. A developer using
Windows locally and a Linux CI runner can share one release graph and lifecycle
while choosing a shell explicitly. Keep these three files in the repository root:
\begin{speccode}
# global.yaml: shared graph, scripts and lifecycle
scripts:
  build: cmake --build build
flow:
  build: [build]
packages:
  app:
    path: app

# linux.yaml: Linux overrides only
$ref: ./global.yaml
shell: ["/bin/sh", "-c"]

# windows.yaml: Windows overrides only
$ref: ./global.yaml
shell: ["powershell.exe", "-NoProfile", "-NonInteractive",
        "-Command"]
\end{speccode}
The platform files inherit the shared configuration and replace only
\code{shell}. References resolve relative to the file containing them.
Overrides replace whole keys, not nested fields: overriding \code{scripts}
would replace that entire map. Individual scripts can instead reference
separate files with \code{\$ref}.

From Windows, preview with \code{dispat --config windows.yaml status};
on the Linux runner, use \code{dispat --config linux.yaml status}.
After its release gates, CI can run \code{dispat --config linux.yaml} to
execute the release. Where a compatible POSIX shell is available, one selector
can choose the preview configuration:
\begin{speccode}
dispat if OS=Windows_NT \
  --then 'dispat --config windows.yaml status' \
  --else 'dispat --config linux.yaml status'
\end{speccode}
\code{OS=Windows\_NT} tests an environment value; it is not host detection.
Here the non-Windows branch assumes a Linux runner. The evaluated releases
execute \code{if} branches through \code{/bin/sh -c} without reading
configuration; the child invocation then loads its platform shell. Thus the
selector needs \code{/bin/sh}, even though the Windows configuration names
PowerShell. A native Windows user can invoke the Windows configuration directly
instead. The example assumes CMake is installed and \code{app/build} is already
configured. It does not establish native Windows execution coverage.

A shared \code{dispat for Debug Release --do "dispat exec build-one"}
can likewise pass each build configuration through \code{DISPAT\_ITEM} to a
named script. A configuration-free literal \code{for} loop also requires
\code{/bin/sh}; platform-specific leaf commands must still be valid for the
chosen shell. These are composition examples, not additional protocol guarantees.

For an agent-assisted release, complementary gates keep evidence attached to
the inputs being published: a lock-free
\code{dispat status --require-release --strict}; an all-package workspace suite
with temporary links and cleanup that runs and is verified under
\code{always()}; and package checks after version/lockfile reconciliation plus
smoke tests of the exact built artefact before publish. These commands inspect
or test inputs; they do not implicitly execute the computed release plan. Every wrapper must
preserve exit status, and these checks neither start a release nor change the
required authorization policy. HOL Guard PR~\cite{holguard} adds opt-in review
when a Dispat release starts, but at its cited open, unmerged revision it is not
an implementation or experimental validation of this whole gate stack; that
integration remains in progress.

\subsection{How to read it}

\Cref{sec:model} is definitional and dense; a reader who wants the results
first can read \cref{sec:guarantees} directly, following the macro names back
to \cref{sec:model} as needed. \Cref{sec:execution} can be read on its own by
anyone interested in what placement and the readiness relation can and cannot
buy; \cref{sec:evaluation} likewise, by anyone who wants the numbers and their
provenance before the mathematics. \Cref{appendix:notation} is the symbol
table. Every definition names the specification section it transcribes, and
every theorem names the guarantee it formalises, so the report can be read
beside its source. \Cref{sec:conformance} maps each result to the conformance
vectors of \SPEC{} Appendix B and to the test in the reference implementation
that exercises it.

\subsection{Licence}

\SPEC{} is licensed GPL-3.0-or-later, and this report reproduces its
definitions and pseudocode; it is therefore distributed under the same licence.
See \code{paper/LICENSE}.
